Considera la funció
\[
f(x)=\frac{6}{x+1} .
\]

\begin{apartats}

\apartat[1,25]{0,75}
Calcula la taxa de variació mitjana de $f$ als intervals $[0,2]$ i $[1,5]$, i interpreta'n el
resultat.

\begin{solucio}
$f(0)=6$, $f(1)=3$, $f(2)=2$ i $f(5)=1$.\\
$\mathrm{TVM}[0,2]=\dfrac{2-6}{2}=-2$ \quad i \quad $\mathrm{TVM}[1,5]=\dfrac{1-3}{4}=-0{,}5$.\\
Les dues són negatives: la funció decreix als dos intervals. A més, decreix molt més de pressa
a prop de $x=0$ que no pas entre $x=1$ i $x=5$.
\end{solucio}

\apartat[1,25]{1}
Calcula $f'(1)$ aplicant la definició de derivada.

\begin{solucio}
\[
f'(1)=\lim_{h\to0}\frac{\frac{6}{2+h}-3}{h}
=\lim_{h\to0}\frac{6-3(2+h)}{h(2+h)}
=\lim_{h\to0}\frac{-3h}{h(2+h)}
=\lim_{h\to0}\frac{-3}{2+h}=-\frac32 .
\]
\end{solucio}

\begin{nomesllarg}
\apartat{0,75}
Calcula $f'(4)$ aplicant la definició i compara'l amb $f'(1)$. Què indica sobre la gràfica?

\begin{solucio}
\[
f'(4)=\lim_{h\to0}\frac{\frac{6}{5+h}-\frac65}{h}
=\lim_{h\to0}\frac{30-6(5+h)}{5h(5+h)}
=\lim_{h\to0}\frac{-6}{5(5+h)}=-\frac{6}{25} .
\]
Totes dues derivades són negatives, però $|f'(1)|=1{,}5$ és molt més gran que
$|f'(4)|=0{,}24$: la gràfica baixa molt més inclinada en $x=1$ que en $x=4$, on ja s'aplana.
\end{solucio}
\end{nomesllarg}

\end{apartats}
